% Options for packages loaded elsewhere
\PassOptionsToPackage{unicode}{hyperref}
\PassOptionsToPackage{hyphens}{url}
\documentclass[
]{article}
\usepackage{xcolor}
\usepackage{amsmath,amssymb}
\setcounter{secnumdepth}{-\maxdimen} % remove section numbering
\usepackage{iftex}
\ifPDFTeX
  \usepackage[T1]{fontenc}
  \usepackage[utf8]{inputenc}
  \usepackage{textcomp} % provide euro and other symbols
\else % if luatex or xetex
  \usepackage{unicode-math} % this also loads fontspec
  \defaultfontfeatures{Scale=MatchLowercase}
  \defaultfontfeatures[\rmfamily]{Ligatures=TeX,Scale=1}
\fi
\usepackage{lmodern}
\ifPDFTeX\else
  % xetex/luatex font selection
\fi
% Use upquote if available, for straight quotes in verbatim environments
\IfFileExists{upquote.sty}{\usepackage{upquote}}{}
\IfFileExists{microtype.sty}{% use microtype if available
  \usepackage[]{microtype}
  \UseMicrotypeSet[protrusion]{basicmath} % disable protrusion for tt fonts
}{}
\makeatletter
\@ifundefined{KOMAClassName}{% if non-KOMA class
  \IfFileExists{parskip.sty}{%
    \usepackage{parskip}
  }{% else
    \setlength{\parindent}{0pt}
    \setlength{\parskip}{6pt plus 2pt minus 1pt}}
}{% if KOMA class
  \KOMAoptions{parskip=half}}
\makeatother
\usepackage{longtable,booktabs,array}
\newcounter{none} % for unnumbered tables
\usepackage{calc} % for calculating minipage widths
% Correct order of tables after \paragraph or \subparagraph
\usepackage{etoolbox}
\makeatletter
\patchcmd\longtable{\par}{\if@noskipsec\mbox{}\fi\par}{}{}
\makeatother
% Allow footnotes in longtable head/foot
\IfFileExists{footnotehyper.sty}{\usepackage{footnotehyper}}{\usepackage{footnote}}
\makesavenoteenv{longtable}
\usepackage{graphicx}
\makeatletter
\newsavebox\pandoc@box
\newcommand*\pandocbounded[1]{% scales image to fit in text height/width
  \sbox\pandoc@box{#1}%
  \Gscale@div\@tempa{\textheight}{\dimexpr\ht\pandoc@box+\dp\pandoc@box\relax}%
  \Gscale@div\@tempb{\linewidth}{\wd\pandoc@box}%
  \ifdim\@tempb\p@<\@tempa\p@\let\@tempa\@tempb\fi% select the smaller of both
  \ifdim\@tempa\p@<\p@\scalebox{\@tempa}{\usebox\pandoc@box}%
  \else\usebox{\pandoc@box}%
  \fi%
}
% Set default figure placement to htbp
\def\fps@figure{htbp}
\makeatother
\setlength{\emergencystretch}{3em} % prevent overfull lines
\providecommand{\tightlist}{%
  \setlength{\itemsep}{0pt}\setlength{\parskip}{0pt}}
% ---------------------------------------------------------------------------
% Unicode -> LaTeX mappings so these documents compile and render correctly
% under xelatex (and pdflatex) without depending on a particular system font.
% Every non-ASCII glyph used in the course text is mapped explicitly; this is
% needed because the documents use mathematical symbols inside running text
% and inside inline-code spans, where unicode-math would otherwise treat them
% as math-only and abort.
% ---------------------------------------------------------------------------
\usepackage{amssymb}
\usepackage{newunicodechar}

% --- Greek (lower case) ---
\newunicodechar{α}{\ensuremath{\alpha}}
\newunicodechar{β}{\ensuremath{\beta}}
\newunicodechar{γ}{\ensuremath{\gamma}}
\newunicodechar{δ}{\ensuremath{\delta}}
\newunicodechar{θ}{\ensuremath{\theta}}
\newunicodechar{κ}{\ensuremath{\kappa}}
\newunicodechar{λ}{\ensuremath{\lambda}}
\newunicodechar{μ}{\ensuremath{\mu}}
\newunicodechar{ν}{\ensuremath{\nu}}
\newunicodechar{π}{\ensuremath{\pi}}
\newunicodechar{ρ}{\ensuremath{\rho}}
\newunicodechar{ς}{\ensuremath{\varsigma}}
\newunicodechar{σ}{\ensuremath{\sigma}}
\newunicodechar{τ}{\ensuremath{\tau}}
\newunicodechar{φ}{\ensuremath{\varphi}}
\newunicodechar{χ}{\ensuremath{\chi}}
% --- Greek (capitals) ---
\newunicodechar{Γ}{\ensuremath{\Gamma}}
\newunicodechar{Δ}{\ensuremath{\Delta}}
\newunicodechar{Λ}{\ensuremath{\Lambda}}
\newunicodechar{Σ}{\ensuremath{\Sigma}}

% --- Math symbols ---
\newunicodechar{ℝ}{\ensuremath{\mathbb{R}}}
\newunicodechar{∂}{\ensuremath{\partial}}
\newunicodechar{∈}{\ensuremath{\in}}
\newunicodechar{∫}{\ensuremath{\int}}
\newunicodechar{√}{\ensuremath{\surd}}
\newunicodechar{≈}{\ensuremath{\approx}}
\newunicodechar{≤}{\ensuremath{\leq}}
\newunicodechar{≥}{\ensuremath{\geq}}
\newunicodechar{±}{\ensuremath{\pm}}
\newunicodechar{×}{\ensuremath{\times}}
\newunicodechar{−}{\ensuremath{-}}
\newunicodechar{·}{\ensuremath{\cdot}}
\newunicodechar{½}{\textonehalf}
\newunicodechar{→}{\ensuremath{\rightarrow}}
\newunicodechar{⟶}{\ensuremath{\longrightarrow}}

% --- Superscripts / modifier letters ---
\newunicodechar{²}{\ensuremath{^{2}}}
\newunicodechar{³}{\ensuremath{^{3}}}
\newunicodechar{⁺}{\ensuremath{^{+}}}
\newunicodechar{⁻}{\ensuremath{^{-}}}
\newunicodechar{ᵀ}{\ensuremath{^{\mathsf{T}}}}
\newunicodechar{ᵏ}{\ensuremath{^{k}}}

% --- Punctuation / misc text symbols ---
\newunicodechar{…}{\textellipsis}
\newunicodechar{§}{\S}
\newunicodechar{⚑}{\textbf{[!]}}

% --- Accented Latin (author names) ---
\newunicodechar{Á}{\'A}
\newunicodechar{Ø}{\O}
\newunicodechar{ä}{\"a}
\newunicodechar{é}{\'e}
\newunicodechar{ó}{\'o}
\newunicodechar{ô}{\^o}
\usepackage{bookmark}
\IfFileExists{xurl.sty}{\usepackage{xurl}}{} % add URL line breaks if available
\urlstyle{same}
\hypersetup{
  pdftitle={Optimal Market Making for Cryptocurrency Options},
  pdfauthor={North Carolina State University · Financial Mathematics},
  hidelinks,
  pdfcreator={LaTeX via pandoc}}

\title{Optimal Market Making for Cryptocurrency Options}
\usepackage{etoolbox}
\makeatletter
\providecommand{\subtitle}[1]{% add subtitle to \maketitle
  \apptocmd{\@title}{\par {\large #1 \par}}{}{}
}
\makeatother
\subtitle{The Avellaneda--Stoikov Framework and its Extension to
Options}
\author{North Carolina State University · Financial Mathematics}
\date{Summer 2026}

\begin{document}
\maketitle

\includegraphics[width=2.6in,height=\textheight,keepaspectratio]{assets/ncsu_logo.png}

\section{Optimal Market Making for Cryptocurrency
Options}\label{optimal-market-making-for-cryptocurrency-options}

\textbf{The Avellaneda--Stoikov Framework and its Extension to Options}

A 10-Week Summer Research Program

North Carolina State University --- Department of Mathematics

\textbf{Industry Mentor:} Dendi Suhubdy --- \emph{Bitwyre}, San
Francisco

\textbf{Guest Practitioner:} Fenni Kang --- Chief Strategy Officer,
\emph{Coincall}

Edition: 2026

\begin{center}\rule{0.5\linewidth}{0.5pt}\end{center}

\subsection{1. Course Overview}\label{course-overview}

This ten-week summer research program prepares students in the North
Carolina State University Financial Mathematics graduate program to
design, implement, and evaluate optimal market making strategies for
cryptocurrency options. The program is organized around the
\textbf{Avellaneda--Stoikov framework}, the foundational
stochastic-control approach to optimal liquidity provision. The central
pedagogical question is how to extend Avellaneda--Stoikov, which was
originally formulated for a single linear instrument, to non-linear
derivatives whose inventory risk is multi-dimensional in the Greeks. A
dedicated advanced week treats extensions of Avellaneda--Stoikov to the
multi-asset, adverse-selection, discrete-inventory, signal-driven, and
robust settings.

Week 1 was an \textbf{introduction week}, delivered last week: it
covered the program goals, the development environment, the Coincall
datasets, and the logistics described in Sections 1.5--1.7. The nine
technical weeks that follow (Weeks 2--10) are the substance of this
syllabus.

The program is mentored by Dendi Suhubdy and includes two guest sessions
from Fenni Kang, Chief Strategy Officer at Coincall, on \textbf{taker
strategies} (Week 9) and on the \textbf{corporate usage of options and
structured products} (Week 10). Students will work with live and
historical Coincall data throughout the program and will produce a
research-grade final report and code repository suitable for further
academic or industry use.

\subsubsection{1.1 Implementation
Languages}\label{implementation-languages}

The program uses a deliberate two-language split, and students should
understand the rationale:

\begin{itemize}
\item
  \textbf{Problem sets are written in Python}, using \textbf{PyTorch} as
  the numerical engine. PyTorch is chosen so that \emph{differentiation
  is handled implicitly by autograd}: students compute option Greeks,
  calibrate surfaces, and fit intensities by writing the forward
  computation only and letting \texttt{.backward()} produce the
  sensitivities. We deliberately do \textbf{not} ask students to
  hand-derive adjoint algorithmic differentiation; the goal is correct,
  readable research code, not AD plumbing.
\item
  \textbf{Fast-computation methods are written in C++.} The
  latency-critical numerical kernels --- characteristic-function
  pricers, implied-volatility inversion, and full-surface revaluation
  --- are implemented in C++ for speed, then exposed to Python through a
  \textbf{\texttt{pybind11}} binding so that the same Python research
  code can call them. Week 7 teaches \texttt{pybind11} explicitly: how
  to write a binding module, build it with CMake, manage the GIL, and
  pass NumPy/PyTorch tensors across the boundary without copying.
\end{itemize}

The result is a realistic hybrid that mirrors industry practice: Python
for research and orchestration, C++ for the hot path, and a clean binder
between them.

\subsubsection{1.2 Learning Objectives}\label{learning-objectives}

Upon successful completion of this program, students will be able to:

\begin{enumerate}
\def\labelenumi{\arabic{enumi}.}
\tightlist
\item
  Formulate a market making problem as a continuous-time stochastic
  optimal control problem and understand the associated
  Hamilton--Jacobi--Bellman (HJB) equation.
\item
  Solve the HJB system for the Avellaneda--Stoikov model in closed form
  and implement the resulting quoting policy as a Python/PyTorch engine
  validated against the analytical solution.
\item
  Calibrate parametric volatility surfaces (SVI, SABR) to live
  cryptocurrency options data in PyTorch and assess their stability
  across time.
\item
  Implement latency-critical pricing, implied-volatility, and Greek
  routines in C++ and call them from Python through a \texttt{pybind11}
  binding.
\item
  Extend the Avellaneda--Stoikov framework to a multi-Greek inventory
  penalty following El Aoud and Abergel, and reduce the resulting HJB
  system to a numerically tractable form.
\item
  Integrate a dynamic hedging policy with the quoting engine and analyze
  the resulting P\&L attribution.
\item
  Build and validate an event-driven backtest that consumes real
  Coincall order-book data and produces interpretable P\&L, risk, and
  Greek-exposure diagnostics.
\item
  Situate the maker (market-making) perspective against \textbf{taker
  strategies}, and understand how options are used by corporates through
  \textbf{structured products} (collars, accumulators/decumulators,
  coupon products).
\item
  Implement at least one advanced extension of Avellaneda--Stoikov and
  communicate the results in writing and orally at a level suitable for
  industry quantitative-research interviews.
\end{enumerate}

\subsubsection{1.3 Prerequisites}\label{prerequisites}

Students are expected to have completed or be concurrently enrolled in
graduate coursework covering: stochastic calculus (Itô integration,
Itô's lemma, Girsanov's theorem), continuous-time stochastic processes
(Brownian motion, Poisson processes, jump-diffusions),
Black--Scholes--Merton option pricing, partial differential equations at
the level of Evans, and probability theory at the level of Durrett.

Programming maturity in \textbf{Python} is required, including comfort
with NumPy and the basics of \textbf{PyTorch} (tensors, autograd).
Familiarity with \textbf{modern C++ (C++17)} and \textbf{CMake} is
advantageous for the fast-computation week but is taught from a working
scaffold; no prior \texttt{pybind11} experience is assumed.

\subsubsection{1.4 Structure}\label{structure}

The program runs for ten weeks. Week 1 was the introduction week
(delivered last week). Each of the nine technical weeks (Weeks 2--10)
comprises two ninety-minute lectures, one ninety-minute recitation, and
approximately fifteen hours of independent reading, programming, and
benchmarking. \textbf{Weeks 5, 7, and 9 are milestone weeks} at which
students deliver a working software component built on the preceding
weeks' material. The final week (Week 10) surveys advanced
Avellaneda--Stoikov extensions and is devoted to a chosen extension
topic and the preparation and presentation of a research-grade final
report.

The schedule is built around the Avellaneda--Stoikov framework as the
unifying worked example: it is introduced in Week 3, placed within the
general continuous-time control template in Week 4, implemented as a
production engine in Week 5, and extended to a multi-Greek options
setting in Week 8. Extensions of Avellaneda--Stoikov are consolidated
into the dedicated advanced week (Week 10) rather than being interleaved
throughout.

\subsubsection{1.5 Assessment}\label{assessment}

Final grades are determined as follows. Weekly problem sets account for
thirty percent of the grade. The three milestone deliverables in Weeks
5, 7, and 9 jointly account for thirty percent. The final report and
presentation account for thirty percent. Active participation in
lecture, recitation, and code review accounts for the remaining ten
percent.

The weekly problem sets are distributed as a separate \textbf{Problem
Sets} booklet; fully worked solutions are distributed (to graders, and
to students after each deadline) as a separate \textbf{Solutions}
booklet.

\subsubsection{1.6 Mentorship Model}\label{mentorship-model}

Students will have access to weekly thirty-minute one-on-one office
hours with the industry mentor, in addition to scheduled lectures and
recitations. Office hours are intended for clarification of difficult
material, code review of milestone deliverables, and discussion of
research directions for the final week.

\subsubsection{1.7 Guest Practitioner
Sessions}\label{guest-practitioner-sessions}

Fenni Kang (Chief Strategy Officer, Coincall) contributes two guest
sessions, integrated as additional lectures within the regular schedule:

\begin{itemize}
\tightlist
\item
  \textbf{Week 9 --- Taker Strategies.} A practitioner's view of
  liquidity \emph{consumption}: when and why takers cross the spread,
  the signals that drive aggressive flow, and how taker behaviour shapes
  the order flow that the market maker must price.
\item
  \textbf{Week 10 --- Corporate Usage of Options and Structured
  Products.} How corporates and treasuries use options: collars to fix a
  price band, accumulators and decumulators for staged
  accumulation/liquidation, and coupon (yield-enhancement) products. The
  session connects the market-maker's book to the structured-product
  flow that generates it.
\end{itemize}

\begin{center}\rule{0.5\linewidth}{0.5pt}\end{center}

\subsection{2. Program Milestones at a
Glance}\label{program-milestones-at-a-glance}

The program is organized around three milestone deliverables. Each
milestone integrates the material from the preceding weeks and produces
a software artifact that is incorporated into the final research
repository.

{\def\LTcaptype{none} % do not increment counter
\begin{longtable}[]{@{}
  >{\raggedright\arraybackslash}p{(\linewidth - 4\tabcolsep) * \real{0.2308}}
  >{\raggedright\arraybackslash}p{(\linewidth - 4\tabcolsep) * \real{0.2692}}
  >{\raggedright\arraybackslash}p{(\linewidth - 4\tabcolsep) * \real{0.5000}}@{}}
\toprule\noalign{}
\begin{minipage}[b]{\linewidth}\raggedright
Week
\end{minipage} & \begin{minipage}[b]{\linewidth}\raggedright
Title
\end{minipage} & \begin{minipage}[b]{\linewidth}\raggedright
Deliverable
\end{minipage} \\
\midrule\noalign{}
\endhead
\bottomrule\noalign{}
\endlastfoot
5 & Linear-asset Avellaneda--Stoikov engine & Python/PyTorch
implementation of the Avellaneda--Stoikov model on a single linear
asset, validated against simulated dynamics and benchmarked on one month
of Coincall BTC-PERP order-book data. \\
7 & Fast pricing and Greeks library & C++ pricing/implied-vol/Greeks
library for European options under Black--Scholes and Heston dynamics,
exposed to Python via \texttt{pybind11}, with PyTorch reference Greeks.
Latency target: full Coincall BTC surface revaluation in under 5
milliseconds. \\
9 & Full options market-making backtest & Event-driven backtest of the
multi-Greek market-making engine on three months of Coincall BTC and ETH
options data, with full P\&L attribution, risk metrics, and
Greek-exposure time series. \\
\end{longtable}
}

\begin{center}\rule{0.5\linewidth}{0.5pt}\end{center}

\subsection{3. Week-by-Week Schedule}\label{week-by-week-schedule}

Each of the following sections gives the learning objectives, lecture
outlines, recitation activity, required reading, and (where applicable)
milestone deliverable for one week of the program. Detailed problem-set
questions live in the separate Problem Sets booklet; each week below
lists the problem set's scope and its implementation language.

\begin{center}\rule{0.5\linewidth}{0.5pt}\end{center}

\subsubsection{Week 1 --- Introduction
(Delivered)}\label{week-1-introduction-delivered}

\emph{Theme: orientation, environment, and data. This week was delivered
last week and is recorded here for completeness.}

Week 1 introduced the program's goals and arc, the maker-versus-taker
framing, and the development environment. Students set up the Python +
PyTorch toolchain, a C++17 toolchain with CMake, and a working
\texttt{pybind11} ``hello world'' binding; received credentials to the
Coincall data store; and signed the data-use agreement. There was no
graded problem set for Week 1.

\begin{center}\rule{0.5\linewidth}{0.5pt}\end{center}

\subsubsection{Week 2 --- Foundations of Market
Microstructure}\label{week-2-foundations-of-market-microstructure}

\emph{Theme: from limit order books to the economic role of the market
maker.}

\textbf{Learning Objectives}

\begin{itemize}
\tightlist
\item
  Reconstruct a limit order book from raw exchange feeds and compute
  standard liquidity metrics.
\item
  Articulate the three classical components of the bid-ask spread and
  identify them in real data.
\item
  Explain the role of the market maker as a continuously quoting
  liquidity provider, and contrast it with directional and taker
  trading.
\end{itemize}

\textbf{Lecture 1 --- Limit Order Book Mechanics.} Order types (market,
limit, IOC, post-only, hidden); price-time priority and the matching
algorithm; L1/L2/L3 order-book views; tick size and lot size; and the
Coincall exchange architecture (contract specifications, fee structure,
market maker tiers).

\textbf{Lecture 2 --- Spread Decomposition and the Role of the Market
Maker.} The bid-ask spread decomposed into adverse selection, inventory,
and order-processing components, via the classical models of Roll
(1984), Glosten--Milgrom (1985), and Kyle (1985); realized spread,
effective spread, and price impact; why crypto microstructure differs
from equities.

\textbf{Recitation.} Reconstruct the order book from a Coincall
tick-level recording in Python: replay one hour of BTC-PERP events and
produce a time series of mid-prices, spreads, and book depth. Output: a
Python class that maintains live order-book state from raw events.

\textbf{Required Reading.} Cartea, Jaikumar, and Penalva (2015), Ch.
1--3; Roll (1984); Glosten and Milgrom (1985).

\textbf{Problem Set 2 (Python).} Order-book reconstruction, the Roll
estimator, the Glosten--Milgrom quote-setter (belief update simulated in
PyTorch), and a realized-spread / price-impact decomposition on Coincall
data. See the Problem Sets booklet.

\begin{center}\rule{0.5\linewidth}{0.5pt}\end{center}

\subsubsection{Week 3 --- The Avellaneda--Stoikov
Model}\label{week-3-the-avellanedastoikov-model}

\emph{Theme: a fully worked example of optimal market making through
stochastic control.}

\textbf{Learning Objectives}

\begin{itemize}
\tightlist
\item
  Set up the Avellaneda--Stoikov problem as a continuous-time stochastic
  control problem with exponential utility.
\item
  Understand the associated HJB equation and its reduction to a system
  of ODEs through the standard ansatz.
\item
  Interpret the closed-form reservation price and optimal half-spread
  economically.
\end{itemize}

\textbf{Lecture 1 --- Setup and Dynamics.} Mid-price as arithmetic
Brownian motion; Poisson order flow with intensity λ(δ) = A·exp(−κδ);
wealth and inventory dynamics; exponential (CARA) utility and why it
makes the problem tractable; the terminal liquidation penalty.

\textbf{Lecture 2 --- HJB Solution.} The dynamic programming principle
and the HJB equation; the separation ansatz V(t,x,S,q) =
−exp(−γ(x+qS))·θ(t,q); reduction to ODEs in θ; the reservation price r =
S − qγσ²(T−t); and the optimal half-spread and its decomposition into
risk-aversion and order-flow components.

\textbf{Recitation.} Walk through the derivation together, then
implement the closed-form quotes in PyTorch and visualize the
reservation price and spread as functions of inventory.

\textbf{Required Reading.} Avellaneda and Stoikov (2008); Cartea,
Jaikumar, and Penalva (2015), Ch. 10.

\textbf{Problem Set 3 (Python/PyTorch).} A PyTorch ODE solver for θ(t,q)
validated against the closed-form quotes; a reusable quote engine; a
one-day simulation comparing optimal vs.~symmetric quoting; a γ
sensitivity sweep; and a price-impact experiment. See the Problem Sets
booklet.

\begin{center}\rule{0.5\linewidth}{0.5pt}\end{center}

\subsubsection{Week 4 --- Stochastic Optimal Control in Continuous
Time}\label{week-4-stochastic-optimal-control-in-continuous-time}

\emph{Theme: the general theory underlying the previous week's worked
example.}

\textbf{Learning Objectives}

\begin{itemize}
\tightlist
\item
  State and apply the dynamic programming principle in a general
  continuous-time setting.
\item
  Understand the HJB equation for a generic control problem and the role
  of the verification theorem.
\item
  Recognize the role of CARA and CRRA utilities in producing tractable
  problems.
\end{itemize}

\textbf{Lecture 1 --- Dynamic Programming and the HJB Equation.}
Controlled diffusions and admissible controls; the dynamic programming
principle; the infinitesimal generator; derivation of the HJB equation;
boundary and terminal conditions.

\textbf{Lecture 2 --- Verification Theorems and Tractable Utilities.}
The verification theorem and smooth-fit conditions; viscosity solutions
as the setting for non-smooth value functions; CARA, CRRA, and HARA
families; why CARA decouples the value function from wealth in market
making.

\textbf{Recitation.} Worked solutions to Merton's portfolio problem
under CARA and CRRA, and an optimal-liquidation problem with linear
impact, identifying the recurring pattern: HJB → ansatz → ODE →
solution.

\textbf{Required Reading.} Pham (2009), Ch. 3--4; Cartea, Jaikumar, and
Penalva (2015), Ch. 6.

\textbf{Problem Set 4 (Python/PyTorch).} PyTorch solvers for the Merton
problem under CRRA and CARA; an Almgren--Chriss optimal-execution
solver; an HJB-residual verification harness; and an obstacle-problem
(variational-inequality) solver illustrating where viscosity solutions
are needed. See the Problem Sets booklet.

\begin{center}\rule{0.5\linewidth}{0.5pt}\end{center}

\subsubsection{Week 5 --- The Avellaneda--Stoikov Engine:
Implementation, Calibration, and
Backtest}\label{week-5-the-avellanedastoikov-engine-implementation-calibration-and-backtest}

\emph{Theme: turning the closed-form solution into a calibrated,
backtested quoting engine. \textbf{Milestone 1.}}

\textbf{Learning Objectives}

\begin{itemize}
\tightlist
\item
  Implement the Avellaneda--Stoikov quotes as a clean, tested
  Python/PyTorch engine validated against the analytical solution.
\item
  Calibrate the order-flow intensity functions from historical Coincall
  fill data by maximum likelihood.
\item
  Build an event-driven backtest harness and benchmark the engine on one
  month of Coincall BTC-PERP data.
\end{itemize}

\textbf{Lecture 1 --- From Closed Form to Code.} Engine architecture;
state management (inventory, cash, mid-price estimation from the L2
book); mapping continuous-time quantities to discrete event arrivals;
calibrating λ(δ) = A·exp(−κδ) from fill data; numerical pitfalls (units,
cash-process sign conventions, inventory indexing).

\textbf{Lecture 2 --- Calibration, Simulation, and Validation.}
Maximum-likelihood calibration of A and κ with goodness-of-fit; the
backtest harness (replay engine, fill simulator, P\&L attribution);
validating the P\&L distribution against theory; reproducibility under a
fixed seed.

\textbf{Recitation.} Code review of student engines in progress; each
student walks through one simulated trading day; group resolution of
common bugs; walk-through of the Milestone 1 acceptance criteria.

\textbf{Required Reading.} Avellaneda and Stoikov (2008) (re-read with
implementation in mind); Guéant (2016), Ch. 4 (single-asset market
making, engineering background).

\textbf{Problem Set 5 (Python/PyTorch).} Constitutes Milestone 1: an
optimized quote path, ML intensity calibration, reproduction of the
paper's inventory/P\&L behaviour, a one-month backtest vs.~a symmetric
baseline, and a \texttt{pytest} suite with reproducible seeds. See the
Problem Sets booklet.

\begin{quote}
\textbf{Milestone Deliverable --- Milestone 1: Linear-Asset
Avellaneda--Stoikov Engine.} A Python/PyTorch package implementing the
Avellaneda--Stoikov model on a single linear asset, with calibration,
simulation, and backtesting infrastructure, and a reproducible one-month
backtest on Coincall BTC-PERP data documented in a short technical
report.

\emph{Datasets:} \texttt{coincall\_btc\_perp\_l2\_2025Q4.parquet},
\texttt{coincall\_btc\_perp\_trades\_2025Q4.parquet},
\texttt{coincall\_intensity\_calibration\_set.parquet}.

\emph{Acceptance criteria:} quote computation vectorized over inventory
in PyTorch; intensity calibration with χ² goodness-of-fit; backtest P\&L
within ±50\% of the mentor's reference; reproducible under a fixed seed;
\texttt{pytest} coverage of the core quote computation; report ≤ 10
pages.
\end{quote}

\begin{center}\rule{0.5\linewidth}{0.5pt}\end{center}

\subsubsection{Week 6 --- Options Pricing, the Volatility Surface, and
Calibration}\label{week-6-options-pricing-the-volatility-surface-and-calibration}

\emph{Theme: options theory and the practical construction of a
volatility surface.}

\textbf{Learning Objectives}

\begin{itemize}
\tightlist
\item
  Recall the Greeks of European options, including the second-order
  Greeks vanna and volga.
\item
  State the no-static-arbitrage constraints on the volatility surface
  and verify them on parametric surfaces.
\item
  Calibrate SVI and SABR parameterizations to live Coincall options data
  with attention to stability.
\end{itemize}

\textbf{Lecture 1 --- Greeks and the Volatility Surface.} Black--Scholes
and the Greek vector; vanna, volga, and volatility risk; the
implied-volatility smile and term structure; calendar and butterfly
arbitrage; features of the crypto options surface (sparse strikes, wide
wings, quote staleness).

\textbf{Lecture 2 --- Parameterization and Calibration.} Raw SVI and the
arbitrage-free SVI/eSSVI formulation; SVI calibration by nonlinear least
squares; the SABR formula and the Hagan expansion; time-stability of
calibrated parameters; diagnostic plots.

\textbf{Recitation.} Live calibration in PyTorch: load a Coincall chain
snapshot, fit SVI and SABR (using autograd for the Jacobians), and
produce the standard diagnostics.

\textbf{Required Reading.} Gatheral (2006), Ch. 1--5; Gatheral and
Jacquier (2014); Hagan et al.~(2002); Hull (2021), Ch. 15, 19--21.

\textbf{Problem Set 6 (Python/PyTorch).} Black--Scholes price and Greeks
with PyTorch autograd (including vanna/volga); SVI with arbitrage-free
constraint checks; SVI and SABR calibration to a Coincall snapshot; and
a one-month daily-calibration stability study with smoothing. See the
Problem Sets booklet.

\begin{center}\rule{0.5\linewidth}{0.5pt}\end{center}

\subsubsection{Week 7 --- Fast Computation of Prices, Greeks, Implied
Volatility, and
Surfaces}\label{week-7-fast-computation-of-prices-greeks-implied-volatility-and-surfaces}

\emph{Theme: the latency budget of an options market maker, and the C++
numerical methods that meet it. \textbf{Milestone 2.}}

\textbf{Learning Objectives}

\begin{itemize}
\tightlist
\item
  Implement and benchmark characteristic-function pricing methods
  (Carr--Madan FFT and Fang--Oosterlee COS) in C++.
\item
  Implement robust, fast implied-volatility inversion and full-surface
  revaluation in C++.
\item
  Expose the C++ kernels to Python with \texttt{pybind11}, and validate
  Greeks against PyTorch autograd.
\end{itemize}

\textbf{Lecture 1 --- Characteristic-Function Pricing and Implied
Volatility in C++.} The characteristic function under Black--Scholes and
Heston; the Carr--Madan FFT method; the Fang--Oosterlee COS method and
its truncation interval; robust implied-vol inversion (Jäckel-style
rational starts plus a safeguarded Newton/Halley step); vectorizing a
whole surface.

\textbf{Lecture 2 --- Binding C++ to Python with \texttt{pybind11}.}
This is the binder lecture. How a \texttt{pybind11} module is structured
and built with CMake; passing NumPy/PyTorch tensors across the boundary
without copying (buffer protocol, \texttt{py::array\_t}); releasing the
GIL around the C++ hot loop; and how to keep a single source of truth
for Greeks by validating the C++ results against PyTorch autograd. We
rely on PyTorch for \emph{reference} Greeks rather than hand-rolling
adjoint AD in C++ --- autograd is the oracle, C++ is the speed.

\textbf{Recitation.} Profile a naive Python pricer, move the kernel to
C++, bind it with \texttt{pybind11}, and re-benchmark --- targeting at
least an order of magnitude speedup and a sub-5ms full-surface
revaluation.

\textbf{Required Reading.} Carr and Madan (1999); Fang and Oosterlee
(2008); Jäckel (2015), \emph{Let's Be Rational}; the \texttt{pybind11}
documentation (overview, NumPy, and GIL sections).

\textbf{Problem Set 7 (C++ with \texttt{pybind11}, called from Python).}
Constitutes Milestone 2: a C++ COS/Carr--Madan pricer, a robust C++
implied-vol solver, a vectorized surface revaluation, all bound to
Python via \texttt{pybind11}, with Greeks validated against PyTorch
autograd and latency reported as a distribution. See the Problem Sets
booklet.

\begin{quote}
\textbf{Milestone Deliverable --- Milestone 2: Fast Pricing and Greeks
Library.} A C++ library that prices European options under
Black--Scholes and Heston via characteristic-function methods, inverts
implied volatility robustly, computes the full Greek vector, and
revalues the entire calibrated Coincall surface within a strict latency
budget --- all exposed to Python through a \texttt{pybind11} module.
Greeks are validated against an independent PyTorch-autograd reference.

\emph{Datasets:}
\texttt{coincall\_btc\_options\_chain\_snapshots\_2025Q4.parquet},
\texttt{coincall\_calibrated\_svi\_surfaces\_2025Q4.parquet},
\texttt{coincall\_reference\_prices\_2025Q4.parquet}.

\emph{Acceptance criteria:} full BTC surface (300--500 contracts)
revalued with full Greek vector in under 5 ms on one CPU core; max
relative pricing error 1e-4 vs.~reference; Greeks within 1e-4 of PyTorch
autograd; clean \texttt{pybind11} API importable from Python; tests over
all pricers, the implied-vol solver, and the surface adapter; latency
reported as a distribution, not just a mean.
\end{quote}

\begin{center}\rule{0.5\linewidth}{0.5pt}\end{center}

\subsubsection{Week 8 --- Options Market Making Theory and Dynamic
Hedging}\label{week-8-options-market-making-theory-and-dynamic-hedging}

\emph{Theme: extending Avellaneda--Stoikov to a multi-Greek inventory
penalty, and integrating a dynamic hedging policy.}

\textbf{Learning Objectives}

\begin{itemize}
\tightlist
\item
  Set up the El Aoud--Abergel framework, in which inventory is a vector
  of Greeks.
\item
  Understand the multi-Greek HJB system and the inventory penalty as a
  quadratic form.
\item
  Understand delta and gamma hedging and quantify the P\&L attribution
  of a hedged options book.
\end{itemize}

\textbf{Lecture 1 --- Multi-Dimensional Inventory and the Multi-Greek
HJB.} The inventory vector q = (Δ, Γ, ν, \ldots) and its dynamics under
quoting; the quadratic inventory penalty and its economic
interpretation; the multi-Greek HJB and separation ansatz; reduction to
a problem on the inventory vector; verification that it reduces to
Avellaneda--Stoikov in the one-dimensional case.

\textbf{Lecture 2 --- Dynamic Hedging and Integration with Quoting.}
Continuous and discrete delta hedging; the
gamma--theta--realized-variance identity; gamma scalping; vega hedging
and skew risk; external vs.~integrated hedging architectures; perpetual
futures as the delta-hedging instrument and funding-rate considerations
on Coincall.

\textbf{Recitation.} Multi-Greek HJB walk-through, then a PyTorch
hedged-book simulation comparing hedge policies and P\&L distributions.

\textbf{Required Reading.} El Aoud and Abergel (2015); Taleb (1997),
selected chapters; Cartea, Jaikumar, and Penalva (2015), Ch. 11.

\textbf{Problem Set 8 (Python/PyTorch).} Multi-Greek inventory
aggregation and dynamics; a PyTorch solver for the El Aoud--Abergel
reduced problem verified to collapse to Avellaneda--Stoikov in 1-D;
penalty-matrix calibration from Coincall Greek covariances; numerical
verification of the gamma--theta--variance identity; and a threshold
delta-hedger integrated with the Week 5 engine. See the Problem Sets
booklet.

\begin{center}\rule{0.5\linewidth}{0.5pt}\end{center}

\subsubsection{Week 9 --- Backtesting, Risk Analysis, and Taker
Strategies}\label{week-9-backtesting-risk-analysis-and-taker-strategies}

\emph{Theme: assembling the full system and validating it on historical
data, with a practitioner's view of taker flow. \textbf{Milestone 3.}}

\textbf{Learning Objectives}

\begin{itemize}
\tightlist
\item
  Implement an event-driven backtest harness suitable for options market
  making.
\item
  Produce a full P\&L attribution (spread, inventory, gamma, vega,
  hedging).
\item
  Analyze risk metrics and identify regimes of out- and
  under-performance.
\item
  Understand taker strategies and how aggressive flow shapes the maker's
  problem.
\end{itemize}

\textbf{Lecture 1 --- Backtesting Architecture and Risk Metrics.}
Event-driven simulation (clock, event queue, deterministic replay); the
fill simulator; state tracking (cash, inventory by contract, aggregated
Greeks, hedge position); pitfalls (look-ahead bias, unrealistic fills,
expiry handling); Sharpe/Sortino/Calmar, drawdown, and P\&L attribution;
the deflated Sharpe ratio and block-bootstrap confidence intervals.

\textbf{Lecture 2 (Guest --- Fenni Kang, Coincall) --- Taker
Strategies.} A practitioner's view of liquidity consumption: when and
why takers cross the spread; the signals that drive aggressive flow
(momentum, liquidation cascades, news); execution styles (sweep,
TWAP/VWAP, iceberg-taking); and what taker behaviour implies for the
toxicity of the order flow a market maker faces.

\textbf{Recitation.} Code review of the full backtest; each student
presents one regime and walks through the P\&L attribution.

\textbf{Required Reading.} Cartea, Jaikumar, and Penalva (2015), Ch.
10--11 (re-read); Bailey and López de Prado (2014); López de Prado
(2018), Ch. 11--13.

\textbf{Problem Set 9 (Python/PyTorch).} Constitutes Milestone 3: an
event-driven backtest with deterministic replay, a calibrated/validated
fill simulator, a full multi-Greek backtest with risk metrics and Greek
time series, the P\&L attribution decomposition, and a worst-drawdown
forensic with a taker-flow analysis informed by the guest lecture. See
the Problem Sets booklet.

\begin{quote}
\textbf{Milestone Deliverable --- Milestone 3: Full Options
Market-Making Backtest.} A complete backtest of the multi-Greek engine
on three months of Coincall BTC and ETH options data, with full P\&L
attribution, risk metrics, and Greek-exposure time series, reproducible
from a configuration file and documented in a technical report (≤ 20
pages) that includes a drawdown forensic.

\emph{Datasets:} \texttt{coincall\_btc\_options\_l2\_2025Q4.parquet},
\texttt{coincall\_btc\_options\_trades\_2025Q4.parquet},
\texttt{coincall\_eth\_options\_l2\_2025Q4.parquet},
\texttt{coincall\_eth\_options\_trades\_2025Q4.parquet},
\texttt{coincall\_btc\_eth\_perp\_l2\_2025Q4.parquet},
\texttt{coincall\_funding\_rates\_2025Q4.parquet}.

\emph{Acceptance criteria:} reproducible from config; attribution sums
to total P\&L within 1 bp; Sharpe reported with block-bootstrap CI;
Greek exposures bounded within configured limits; report includes a
forensic of at least one drawdown event.
\end{quote}

\begin{center}\rule{0.5\linewidth}{0.5pt}\end{center}

\subsubsection{Week 10 --- Advanced Avellaneda--Stoikov, Structured
Products, and Final
Presentations}\label{week-10-advanced-avellanedastoikov-structured-products-and-final-presentations}

\emph{Theme: advanced extensions of Avellaneda--Stoikov, the corporate
structured-products view, one chosen extension, and the communication of
results.}

\textbf{Learning Objectives}

\begin{itemize}
\tightlist
\item
  Survey advanced extensions of Avellaneda--Stoikov (multi-asset,
  adverse selection, discrete inventory, signal-driven, robust) and
  implement one.
\item
  Understand how corporates use options through structured products.
\item
  Write and present a research-grade report.
\end{itemize}

\textbf{Lecture 1 --- Advanced Avellaneda--Stoikov Extensions.}
Multi-asset market making across BTC and ETH options under a correlated
covariance matrix; adverse-selection-aware quoting and skew trading from
calibrated SVI parameters; signal-driven quoting (machine-learned
order-flow prediction); discrete-inventory formulations; and robust
optimization with a worst-case inventory penalty over a confidence set.
Students choose one extension by end of day.

\textbf{Lecture 2 (Guest --- Fenni Kang, Coincall) --- Corporate Usage
of Options and Structured Products.} How corporates and treasuries use
options: \textbf{collars} to fix a price band on a holding;
\textbf{accumulators / decumulators} for staged accumulation or
liquidation; and \textbf{coupon / yield-enhancement products}. The
session links these structured-product flows to the market-maker's book
--- how a desk warehouses and hedges the resulting Greek exposure ---
and closes with guidance on research-report structure and the oral
presentation.

\textbf{Recitation.} Final dry-run presentations to the cohort, with
feedback focused on the questions an industry audience would ask.

\textbf{Required Reading.} Guéant (2016), the chapter most relevant to
the chosen extension; El Aoud and Abergel (2015) (re-read with the
extension in mind); one additional paper relevant to the chosen
extension; for the guest session, a short structured-products primer
distributed in advance.

\textbf{Problem Set 10 (Python/PyTorch).} A one-page extension proposal;
a PyTorch implementation of the chosen advanced Avellaneda--Stoikov
extension on top of the milestone infrastructure; a quantitative
comparison against the Week 9 baseline; the final report (≤ 30 pages);
and the final twenty-minute presentation. See the Problem Sets booklet.

\begin{center}\rule{0.5\linewidth}{0.5pt}\end{center}

\subsection{4. Comprehensive Reference
List}\label{comprehensive-reference-list}

\subsubsection{4.1 Primary References}\label{primary-references}

Avellaneda, M., and Stoikov, S. (2008). High-frequency trading in a
limit order book. \emph{Quantitative Finance}, 8(3), 217--224.

El Aoud, S., and Abergel, F. (2015). A stochastic control approach to
option market making. \emph{Market Microstructure and Liquidity}, 1(1),
1550006.

Guéant, O. (2016). \emph{The Financial Mathematics of Market Liquidity:
From Optimal Execution to Market Making}. Chapman \& Hall / CRC.

Cartea, Á., Jaikumar, S., and Penalva, J. (2015). \emph{Algorithmic and
High-Frequency Trading}. Cambridge University Press.

\subsubsection{4.2 Volatility Surface and Options
Pricing}\label{volatility-surface-and-options-pricing}

Gatheral, J. (2006). \emph{The Volatility Surface: A Practitioner's
Guide}. Wiley Finance.

Gatheral, J., and Jacquier, A. (2014). Arbitrage-free SVI volatility
surfaces. \emph{Quantitative Finance}, 14(1), 59--71.

Hagan, P. S., Kumar, D., Lesniewski, A. S., and Woodward, D. E. (2002).
Managing smile risk. \emph{Wilmott Magazine}, September, 84--108.

Jäckel, P. (2015). Let's be rational. \emph{Wilmott Magazine}.

Hull, J. C. (2021). \emph{Options, Futures, and Other Derivatives} (11th
ed.). Pearson.

Taleb, N. N. (1997). \emph{Dynamic Hedging: Managing Vanilla and Exotic
Options}. Wiley.

\subsubsection{4.3 Fast Computation
Methods}\label{fast-computation-methods}

Carr, P., and Madan, D. (1999). Option valuation using the fast Fourier
transform. \emph{Journal of Computational Finance}, 2(4), 61--73.

Fang, F., and Oosterlee, C. W. (2008). A novel pricing method for
European options based on Fourier-cosine series expansions. \emph{SIAM
Journal on Scientific Computing}, 31(2), 826--848.

Glasserman, P. (2003). \emph{Monte Carlo Methods in Financial
Engineering}. Springer.

\subsubsection{4.4 Stochastic Optimal
Control}\label{stochastic-optimal-control}

Pham, H. (2009). \emph{Continuous-time Stochastic Control and
Optimization with Financial Applications}. Springer.

Fleming, W. H., and Soner, H. M. (2006). \emph{Controlled Markov
Processes and Viscosity Solutions} (2nd ed.). Springer.

Øksendal, B. (2003). \emph{Stochastic Differential Equations} (6th ed.).
Springer.

\subsubsection{4.5 Market Microstructure}\label{market-microstructure}

Roll, R. (1984). A simple implicit measure of the effective bid-ask
spread. \emph{Journal of Finance}, 39(4), 1127--1139.

Glosten, L. R., and Milgrom, P. R. (1985). Bid, ask and transaction
prices in a specialist market with heterogeneously informed traders.
\emph{Journal of Financial Economics}, 14(1), 71--100.

Kyle, A. S. (1985). Continuous auctions and insider trading.
\emph{Econometrica}, 53(6), 1315--1335.

\subsubsection{4.6 Software and Tooling}\label{software-and-tooling}

Jakob, W., Rhinelander, J., and Moldovan, D. (2017). \emph{pybind11 ---
Seamless operability between C++11 and Python}.
https://github.com/pybind/pybind11

Paszke, A., et al.~(2019). PyTorch: an imperative style,
high-performance deep learning library. \emph{NeurIPS}.

\subsubsection{4.7 Supplementary}\label{supplementary}

Bailey, D. H., and López de Prado, M. (2014). The deflated Sharpe ratio.
\emph{Journal of Portfolio Management}, 40(5), 94--107.

López de Prado, M. (2018). \emph{Advances in Financial Machine
Learning}. Wiley.

Almgren, R., and Chriss, N. (2001). Optimal execution of portfolio
transactions. \emph{Journal of Risk}, 3(2), 5--40.

\begin{center}\rule{0.5\linewidth}{0.5pt}\end{center}

\subsection{Appendix A. Coincall Dataset
Catalog}\label{appendix-a.-coincall-dataset-catalog}

All datasets are delivered as Parquet files on a secure object store
under a non-disclosure agreement. The naming convention is
\texttt{coincall\_\{instrument\}\_\{data\_type\}\_\{period\}.parquet}.

\begin{itemize}
\tightlist
\item
  \textbf{Order book data} --- L2 snapshots (top ten levels per side),
  at 1-second resolution for options and 100 ms for perpetual futures,
  each row carrying bid/ask prices and sizes per level, a strictly
  increasing sequence number, and an exchange timestamp. L2 is
  sufficient for both intensity calibration and the simulation
  environment.
\item
  \textbf{Trade data} --- the full trade tape (contract id, price, size,
  aggressor side, trade id); options trades additionally carry the mark
  implied volatility at trade time, used for surface calibration in Week
  6 and the fill simulator in Week 9.
\item
  \textbf{Surface data} --- calibrated SVI parameters at five-minute
  frequency, used as a baseline for the fast revaluation pipeline in
  Week 7.
\item
  \textbf{Auxiliary data} --- hourly perpetual-futures funding rates
  (for hedging cost accounting in Week 9) and a static snapshot of
  contract specifications.
\end{itemize}

Students receive credentials to a secure S3-compatible endpoint at the
start of the program. Data may be downloaded locally for analysis but
may not be redistributed. Coincall retains the right to update the
dataset catalog during the program with at least one week of notice.

\end{document}
